\section{Formulation}
\label{sec:formulation}

\subsection{The scattering problem}

Let $\Omega\subset\mathbb{R}^3$ be a bounded perfect electric conductor with surface $\Gamma$ and outward unit normal $\bn$. We work in the frequency domain with time dependence $\rme^{-\rmi\omega t}$ and wavenumber $k$. The conductor is illuminated by the plane wave
\begin{equation}
  \Einc(\bx) = \bp\,\rme^{\rmi k\,\bd\cdot\bx}, \qquad |\bd| = 1, \quad \bp\cdot\bd = 0,
  \label{eq:planewave}
\end{equation}
with propagation direction $\bd$ and polarisation $\bp$. We seek the scattered field $\bE$ in the exterior of $\Omega$. It satisfies \cite{colton2013,kirsch2015}
\begin{align}
  \nabla\times\nabla\times\bE - k^2\bE &= 0 && \text{outside } \Omega, \label{eq:maxwell}\\
  \bn\times\bigl(\bE + \Einc\bigr) &= 0 && \text{on } \Gamma, \label{eq:pec}
\end{align}
together with the Silver--M\"uller radiation condition. Only the tangential part of the electric field is prescribed on a conductor. The normal part of the scattered field is not data, and the method must not impose it.

Far from the body the scattered field is an outgoing spherical wave,
\begin{equation}
  \bE(\bx) = \frac{\rme^{\rmi k r}}{r}\,\bF(\hat{\bx}) + O(r^{-2}), \qquad r = |\bx|, \quad \hat{\bx} = \bx/r,
\end{equation}
and $\bF$ is the far-field pattern. The radar cross-section in the direction $\hat\bx$ is $4\pi|\bF(\hat\bx)|^2$.

\subsection{Fundamental solutions}

Let $\Phi(\bx,\by) = \rme^{\rmi k r}/(4\pi r)$ with $r = |\bx-\by|$ be the fundamental solution of the Helmholtz equation, and let $\rhat = (\bx-\by)/r$. The field of an electric dipole at $\by$ with moment $\ba\in\mathbb{C}^3$ is $\bG(\bx,\by)\,\ba$, where
\begin{equation}
  \bG(\bx,\by) = \Bigl(I + \frac{1}{k^2}\nabla\nabla\Bigr)\Phi
  = \Phi\,\Bigl[\Bigl(1+\frac{\rmi}{kr}-\frac{1}{(kr)^2}\Bigr)I-\Bigl(1+\frac{3\rmi}{kr}-\frac{3}{(kr)^2}\Bigr)\rhat\rhat^{\top}\Bigr]
  \label{eq:dyadic}
\end{equation}
is the dyadic Green's function. Each column of $\bG$ satisfies \eqref{eq:maxwell} away from $\by$ and the radiation condition. A second family of solutions is obtained by taking the curl. The field of a magnetic dipole with moment $\bm$ is
\begin{equation}
  \frac{1}{k}\,\nabla_{\bx}\Phi(\bx,\by)\times\bm
  = \Bigl(\rmi-\frac{1}{kr}\Bigr)\Phi\;\rhat\times\bm .
  \label{eq:magnetic}
\end{equation}
For a smooth body the electric dipoles alone are enough. We shall see in \autoref{sec:skeleton} why bodies with edges need both.

\subsection{The method}

We choose source points $\by_1,\dots,\by_N$ inside $\Omega$ and approximate the scattered field by
\begin{equation}
  \bE_N(\bx) = \sum_{j=1}^{N}\Bigl[\bG(\bx,\by_j)\,\ba_j + \frac{1}{k}\nabla_{\bx}\Phi(\bx,\by_j)\times\bm_j\Bigr].
  \label{eq:mfs}
\end{equation}
The moments $\ba_j$ and $\bm_j$ are the unknowns. Where a layout uses only electric dipoles, or only some components, the others are set to zero. The approximation satisfies \eqref{eq:maxwell} and the radiation condition for any choice of the moments. The moments are determined by the boundary condition \eqref{eq:pec}. We choose collocation points $\bx_1,\dots,\bx_M$ on $\Gamma$, with two orthonormal tangent vectors $\bt_i^{1},\bt_i^{2}$ and a positive weight $w_i$ at each point, and minimise
\begin{equation}
  \sum_{i=1}^{M} w_i \sum_{s=1}^{2}\bigl|\bt_i^{s}\cdot\bigl(\bE_N(\bx_i)+\Einc(\bx_i)\bigr)\bigr|^2 .
  \label{eq:leastsquares}
\end{equation}
This is a linear least-squares problem $\min\|A\mathbf{c}-\mathbf{b}\|_2$ with $2M$ equations. We always take more equations than unknowns. The weight $w_i$ is the surface area that the point represents, so that \eqref{eq:leastsquares} approximates the $L^2(\Gamma)$ norm of the tangential residual. For the cube the weights also contain a power of the distance to the nearest edge (\autoref{sec:cube}).

We scale the columns of $A$ to unit norm. With $n$ unknowns the solution costs $O(Mn^2)$ operations, and we use one of two methods. On the sphere and the torus the matrices are ill conditioned, as is usual for the MFS. There we use orthogonal factorisations \cite{golub2013}. For the sphere this is Householder QR, where the factor $Q$ is never formed and the reflectors are applied to the right-hand side. For the small systems of the torus it is the singular value decomposition. For the layouts of the cube and the almond the scaled matrices are well conditioned, with condition numbers below $10^{6}$. There we use the corrected semi-normal equations \cite{bjorck1987}: a Cholesky factorisation of $A^{\mathrm{H}}A$ followed by two steps of refinement with the residual. This costs about half as much as the QR factorisation and gave the same fields to fourteen digits. All systems in this paper are dense and have at most a few thousand unknowns each once the symmetry of the body has been used.

\subsection{How we measure the error}
\label{sec:errors}

We use four measures.

\emph{An exact solution.} For the sphere the Mie series \cite{mie1908,bohren1998} gives the scattered field to machine precision. We report the relative error
\begin{equation}
  \varepsilon = \frac{\bigl(\sum_{l}|\bE_N(\mathbf{z}_l)-\bE(\mathbf{z}_l)|^2\bigr)^{1/2}}{\bigl(\sum_{l}|\bE(\mathbf{z}_l)|^2\bigr)^{1/2}}
  \label{eq:fielderror}
\end{equation}
at 400 points $\mathbf{z}_l$ on a sphere that encloses the body. For the other shapes we use \eqref{eq:fielderror} with a reference field in place of $\bE$, and we always say which.

\emph{The boundary residual.} The square root of \eqref{eq:leastsquares}, divided by the same expression for $\Einc$ alone and evaluated at points that were not used in the fit.

\emph{The optical theorem.} A perfect conductor absorbs no energy, so the extinction and scattering cross-sections
\begin{equation}
  \sext = \frac{4\pi}{k}\,\operatorname{Im}\bigl(\bar\bp\cdot\bF(\bd)\bigr), \qquad
  \ssca = \int_{S^2}|\bF(\hat\bx)|^2\,\mathrm{d}s(\hat\bx)
  \label{eq:optical}
\end{equation}
are equal for the exact solution. For a collocation method they are not equal automatically, and the mismatch $|\sext-\ssca|/\ssca$ is an error indicator that needs no reference solution. On the cube it exposed the fits that were under-sampled at the edges, in which the residual at the collocation points is small and the field is wrong (\autoref{sec:cube}). It is an indicator and not a bound. One fit with too coarse a grid in the interior of the faces had an error of $2.4\%$ and a mismatch of $10^{-5}$.

\emph{A second solver.} We use the boundary element library Bempp-cl \cite{betcke2021,scroggs2017}. It solves the electric field integral equation with Rao--Wilton--Glisson elements \cite{rao1982} on a triangular mesh. We assemble the dense matrix with the OpenCL backend of the library in double precision and solve by LU factorisation. The boundary element solution serves both as an independent reference and as the yardstick for the computing time. All timings in this paper were measured on one machine with four cores and include assembly, solution and the evaluation of the field at the test points. The details are in \ref{sec:bempp}.
